%% Nigerian Digital Travel Registry — Policy Vision and Framework Paper
%% Prepared for submission to the National Information Technology Development Agency (NITDA)
%% Standalone version — not conference-specific
%% Author: Olajide Adewale Olaniyan

\documentclass[11pt]{article}
\usepackage[a4paper, margin=0.9in]{geometry}
\usepackage{titlesec}
\usepackage{booktabs}
\usepackage{tabularx}
\usepackage{xcolor}
\usepackage{enumitem}
\usepackage{tgheros}
\usepackage[T1]{fontenc}
\usepackage{graphicx}
\usepackage{array}
\usepackage[hidelinks]{hyperref}
\usepackage{fancyhdr}
\renewcommand{\familydefault}{\sfdefault}

%% Color definitions
\definecolor{primary-blue}{RGB}{0,51,102}
\definecolor{light-grey}{gray}{0.85}
\definecolor{dark-grey}{gray}{0.4}
\definecolor{accent-green}{RGB}{34,139,34}
\definecolor{dev-orange}{RGB}{204,85,0}

%% Section formatting with colored rules
\titleformat{\section}{
    \color{primary-blue}\bfseries\Large\raggedright
}{}{0em}{}[\color{light-grey}{\titlerule[2pt]}\vspace{-8pt}]

\titleformat{\subsection}{
    \color{primary-blue}\bfseries\large\raggedright
}{}{0em}{}[\vspace{-4pt}]

\titleformat{\subsubsection}{
    \color{primary-blue}\bfseries\normalsize\raggedright
}{}{0em}{}

%% Paragraph spacing
\setlength{\parindent}{0pt}
\setlength{\parskip}{5pt}

%% Hyperlink colours
\hypersetup{
    colorlinks=true,
    linkcolor=primary-blue,
    urlcolor=primary-blue,
    citecolor=dark-grey
}

%% Header / footer
\pagestyle{fancy}
\fancyhf{}
\renewcommand{\headrulewidth}{0.4pt}
\fancyhead[L]{\small\color{dark-grey}Nigerian Digital Travel Registry (NDTR)}
\fancyhead[R]{\small\color{dark-grey}Policy Vision \& Framework Paper}
\fancyfoot[C]{\small\color{dark-grey}\thepage}

%% ======================================================
%%  DOCUMENT BODY
%% ======================================================
\begin{document}
\thispagestyle{empty}

%% ----- Cover page -----
\vspace*{1.2cm}
\begin{center}
    {\color{primary-blue}\rule{\textwidth}{3pt}}\\[14pt]
    {\color{primary-blue}\LARGE\bfseries From Collective Punishment to Evidence-Based Trust}\\[10pt]
    {\color{primary-blue}\large\bfseries A Policy Vision and Framework for Nigeria's Digital Travel Registry\\
    as a Diplomatic Data Asset}\\[12pt]
    {\color{primary-blue}\rule{\textwidth}{1pt}}\\[18pt]

    \noindent
    \begin{tabularx}{\textwidth}{@{}lX@{}}
        \color{dark-grey}\small\textbf{Document type}  & \color{dark-grey}\small Policy Vision and Framework Paper \\[4pt]
        \color{dark-grey}\small\textbf{Document status} & \color{dark-grey}\small Working Draft — For Review and Engagement \\[4pt]
        \color{dark-grey}\small\textbf{Prepared by}    & \color{dark-grey}\small Olajide Adewale Olaniyan \\
                                                        & \color{dark-grey}\small IBM Ireland, Dublin, Ireland \\
                                                        & \color{dark-grey}\small \href{mailto:jidemobell@gmail.com}{jidemobell@gmail.com} \\[4pt]
        \color{dark-grey}\small\textbf{Prepared for}   & \color{dark-grey}\small National Information Technology Development Agency (NITDA) \\[4pt]
        \color{dark-grey}\small\textbf{Date}           & \color{dark-grey}\small April 2026 \\[4pt]
        \color{dark-grey}\small\textbf{Version}        & \color{dark-grey}\small 1.0 \\
    \end{tabularx}

    \vspace{1.2cm}
    \noindent\colorbox{primary-blue!7}{%
    \begin{minipage}{\dimexpr\textwidth-2\fboxsep}
    \vspace{6pt}\centering
    \small\color{dark-grey}
    This document is submitted in connection with NITDA's mandate on digital governance
    and Nigeria's national digital transformation agenda.
    It is intended to inform policy deliberation, feasibility assessment, and potential
    pilot commissioning of a Nigerian Digital Travel Registry.
    \vspace{6pt}
    \end{minipage}}
\end{center}
\clearpage

%% ----- Abstract box -----
\noindent\colorbox{primary-blue!7}{%
\begin{minipage}{\dimexpr\textwidth-2\fboxsep}
\vspace{6pt}
\textbf{\color{primary-blue}Abstract.}\quad
Nigeria's passport remains among the most visa-restricted in the world, a status that
disproportionately penalises the overwhelming majority of Nigerian citizens who travel
lawfully for education, business, and family purposes.
This document argues that the root cause of this diplomatic vulnerability is not a
disproportionate rate of travel-related misconduct, but rather a systemic absence of
verifiable, aggregate travel behaviour data that could demonstrate compliance at scale.
We propose the Nigerian Digital Travel Registry (NDTR)---a mandatory digital travel
declaration framework modelled on the Philippines' eTravel system---as both a governance
modernisation tool and a diplomatic data asset.
The NDTR framework is examined across four functional dimensions:
(i) generating compliance evidence to counter blanket visa restrictions;
(ii) enabling proactive security monitoring through anomaly detection in travel patterns;
(iii) strengthening Nigeria's position in bilateral visa negotiations through
evidence-based advocacy; and
(iv) modernising Nigerian border management in alignment with ICAO, ECOWAS,
and African Union digital governance standards.
Drawing on a comparative analysis of the Philippines eTravel system---whose
data-driven approach contributed to Canada granting visa-free access to Filipino
citizens---this document develops a phased implementation roadmap that integrates
Nigeria's National Identity Number (NIN), the Nigeria Immigration Service (NIS),
and the National Data Protection Commission (NDPC) into a coherent digital migration
governance architecture.
This document presents a policy vision and implementation framework; empirical
validation is proposed as future work following a pilot deployment.
\vspace{4pt}

\textbf{\color{primary-blue}Keywords:}\quad
digital travel registry \textbf{\textperiodcentered} migration governance
\textbf{\textperiodcentered} diplomatic data \textbf{\textperiodcentered} Nigeria
\textbf{\textperiodcentered} visa restrictions \textbf{\textperiodcentered} eTravel
\textbf{\textperiodcentered} NIN \textbf{\textperiodcentered} border management
\textbf{\textperiodcentered} data-driven governance \textbf{\textperiodcentered} NITDA
\vspace{6pt}
\end{minipage}}

\vspace{14pt}

%% ----- Table of Contents -----
{\color{primary-blue}\bfseries\large Contents}\\[-4pt]
{\color{light-grey}\rule{\textwidth}{2pt}}
\vspace{4pt}
\begin{enumerate}[label=\textbf{\arabic*.}, leftmargin=*, itemsep=2pt]
    \item Executive Summary
    \item Introduction
    \item Background: The Data Absence Penalty
    \item Related Work: Digital Travel Declaration Systems
    \item The NDTR Framework
    \item Proactive Security Monitoring and Anomaly Detection
    \item Diplomatic Applications of NDTR Data
    \item Border Management Modernisation
    \item Implementation Roadmap
    \item Discussion
    \item Policy Recommendations
    \item Conclusion
\end{enumerate}
\vspace{10pt}

%% ----------------------------------------------------------
\section*{Executive Summary}
%% ----------------------------------------------------------
\addcontentsline{toc}{section}{Executive Summary}

Nigeria's diplomatic standing in the global migration system is structurally disadvantaged by
a single, correctable governance gap: the absence of a national digital travel declaration
system. Nigerian citizens bear collective reputational and visa access costs attributable to
a small minority of travel-related misconduct, yet the Nigerian government currently lacks
the data infrastructure required to demonstrate---with evidence---that this minority is
statistically marginal.

This document proposes the \textbf{Nigerian Digital Travel Registry (NDTR)}: a mandatory,
NIN-anchored digital pre-departure and arrival declaration system for all Nigerian citizens
travelling internationally. The NDTR is modelled on the Philippines' eTravel system, which
has been credited as a contributing factor in Canada granting visa-free access to Philippine
passport holders in 2024. The NDTR would generate the national-level compliance dataset
that Nigeria currently lacks, converting it into a usable instrument for bilateral
diplomatic negotiations.

\textbf{The NDTR serves four interconnected policy objectives:}
\begin{itemize}[leftmargin=*]
    \item \textbf{Compliance evidence:} Produces verifiable, government-generated data
    demonstrating the travel behaviour of Nigerian citizens at scale, for use in bilateral
    visa negotiations and ECOWAS free movement dialogues.
    \item \textbf{Security monitoring:} Enables proactive, pattern-based anomaly detection
    in outbound travel, replacing the current entirely reactive posture.
    \item \textbf{Asylum integrity:} Provides verifiable travel history data that can
    strengthen the integrity of asylum adjudication processes in destination countries.
    \item \textbf{Border modernisation:} Aligns Nigerian border management with ICAO Doc~9303
    standards, ECOWAS digital free movement objectives, and the African Union digital
    passport infrastructure programme.
\end{itemize}

\textbf{Key recommendations for NITDA:}
\begin{enumerate}[leftmargin=*]
    \item Establish an interagency NDTR Steering Committee chaired by NITDA, with
    NIS, NIMC, NDPC, and the Ministry of Foreign Affairs as core members.
    \item Commission a formal feasibility study and legal gap analysis to identify
    the legislative amendments required.
    \item Initiate a voluntary pilot deployment at Lagos International Airport within
    12 months, integrated with the existing NIN infrastructure.
    \item Publish the first national travel compliance report within 24 months of
    full launch, for use in visa liberalisation dialogues.
    \item Engage bilaterally with at least two destination countries---United Kingdom
    and a Schengen representative---using NDTR data within 36 months.
\end{enumerate}

The Philippines case demonstrates that this approach works.
Nigeria has the institutional infrastructure, the NIN identity backbone, and
the policy impetus to replicate it.
The question is not whether the NDTR is feasible; it is whether Nigeria will
act before the window of diplomatic leverage closes further.

\vspace{10pt}

%% ----------------------------------------------------------
\section{Introduction}
%% ----------------------------------------------------------

The Nigerian passport currently ranks among the least powerful in the world,
granting visa-free or visa-on-arrival access to fewer than 45 destinations
as of 2025~\cite{henley2025}. This places Nigerian citizens in a position of
systematic travel disadvantage compared to peers in countries of comparable
economic development and, in some cases, countries with documented rates of
irregular migration that rival or exceed Nigeria's. The Philippines,
Colombia, and Ghana---all confronting similar immigration challenges at
various points in their recent histories---have progressed toward greater
passport freedom in part by investing in governance systems that generate
verifiable data on the travel behaviour of their citizens.

Nigeria, by contrast, has yet to implement a comprehensive digital travel
declaration system. The consequence is structural: Nigerian immigration
authorities negotiate at the bilateral level without the evidentiary tools
needed to differentiate the travel behaviour of the vast majority of
lawful travellers from a statistically small minority responsible for
immigration violations.
The result is what this document terms a \textit{data absence penalty}---a
measurable degradation in diplomatic standing attributable not to
demonstrated national misconduct, but to the inability to prove its
absence.

This document proposes the Nigerian Digital Travel Registry (NDTR) as the
policy and technical intervention designed to close this gap. The NDTR is
a mandatory, digital, pre-departure and arrival declaration system for all
Nigerian citizens travelling internationally. The framework is modelled on
the Philippines' eTravel system, which the author personally used during
a visit to the Philippines in January 2026, providing firsthand evidence of
the system's usability and operational design. The NDTR goes further, integrating
Nigeria's National Identity Number (NIN) infrastructure to create a persistent,
privacy-preserving travel behaviour profile that can be aggregated at a national
level for diplomatic and security purposes.

The motivation for this work is grounded in a specific policy problem:
the majority of Nigerians who travel internationally do so lawfully,
yet they bear the costs---in visa refusals, heightened scrutiny, and
reputational stigma---of a small minority. Available migration data suggests
that the proportion of Nigerian travellers involved in documented violations
represents a statistically small share of total annual outbound
travel~\cite{iom2023,unodc2022}, yet visa restriction regimes are applied
nationally without this distinction. The NDTR is the instrument by which
this statistical reality can be made visible and actionable.

This document is structured as follows:
Section~\ref{sec:background} situates the problem within the global
migration governance context.
Section~\ref{sec:related} reviews comparable digital travel declaration
systems internationally.
Section~\ref{sec:framework} presents the NDTR architecture and design
principles.
Section~\ref{sec:security} examines the system's proactive security
monitoring capabilities.
Section~\ref{sec:diplomatic} develops the diplomatic applications of NDTR data.
Section~\ref{sec:border} addresses border management modernisation and
international standards alignment.
Section~\ref{sec:roadmap} proposes a phased implementation roadmap.
Section~\ref{sec:recommendations} sets out concrete policy recommendations
for NITDA and partner agencies.
Section~\ref{sec:conclusion} concludes.

%% ----------------------------------------------------------
\section{Background: The Data Absence Penalty}
\label{sec:background}
%% ----------------------------------------------------------

\subsection{Nigeria's Passport Restriction Problem}

Nigeria's Henley Passport Index ranking has remained in the bottom
quartile globally for over a decade~\cite{henley2025}. This is not a
static condition but a cascading one: as visa restrictions accumulate,
they reinforce perceptions that justify further restrictions in a
self-amplifying reputational cycle. The diplomatic cost is significant.
Nigerian professionals, academics, and businesspeople routinely face
visa refusals, processing delays, and travel bans that their peers
from other developing nations do not, regardless of individual
travel history or socioeconomic standing.

Crucially, the restrictions are not uniformly evidence-based.
Several countries that impose the most stringent restrictions on
Nigerian passports---including the Philippines, Colombia, and EU
Schengen member states---do so on the basis of aggregate immigration
enforcement data that does not distinguish between the rare offender
and the frequent lawful traveller. The absence of countervailing
Nigerian government data leaves destination-country immigration
authorities with no alternative evidentiary basis for policy
differentiation.

\subsection{The Collective Punishment Problem}

The term \textit{collective punishment} in the immigration policy context
refers to the application of restrictive migration controls at a
population level, based on the actions of a statistically marginal
subset of that population~\cite{castles2013}. This represents a classic
case of \textit{base rate neglect}~\cite{kahneman2011}: policymakers
respond to high-visibility incidents---a drug courier, a fraudulent
asylum claim, an immigration overstay---without reference to the
statistical base from which those incidents are drawn.

If, conservatively, 5 million Nigerian international trips occur
annually and a small fraction result in documented violations, a
visa restriction applied nationally punishes the overwhelming majority
to control a marginal minority. More critically, it does not effectively
control that minority either---it merely inconveniences lawful travellers
at the point of application while motivated actors find alternative routes.
The practical effect is to harm the legitimate traveller while providing
minimal security gain.

The Nigerian government's diplomatic response to this situation has
historically been rhetorical rather than evidentiary. Without a
national travel data system, Nigeria cannot present the quantitative
case that would shift bilateral conversations from perception to fact.

\subsection{The Asylum Fraud Dimension}

A related but distinct problem is the pattern of asylum claims lodged
by Nigerian nationals in European and North American jurisdictions
that do not reflect genuine persecution risk. Immigration tribunals in
Ireland, the United Kingdom, and Germany have noted a pattern of
Nigerian asylum applicants who have demonstrable employment histories,
prior international travel records, and no documented nexus to regions
affected by insecurity in Nigeria---yet claim refugee status on the
basis of generalised fear~\cite{eurostat2024,irishnat2023}.

The absence of verifiable, longitudinal travel and socioeconomic data
from the Nigerian government means that receiving-country adjudicators
must make credibility determinations without access to the kind of
``digital footprint'' profile that would make fabrication substantially
harder. An individual with a documented history of multiple international
business trips across several years, registered through the NDTR with NIN
linkage, presents a very different profile to an asylum adjudicator
than a claimant whose travel record is entirely opaque.
This cross-border data sharing capability---with appropriate bilateral
agreements and privacy safeguards---would simultaneously deter
fraudulent claims and protect the credibility of genuine ones.

%% ----------------------------------------------------------
\section{Related Work: Digital Travel Declaration Systems}
\label{sec:related}
%% ----------------------------------------------------------

A range of digital travel authorisation and declaration systems have
been deployed globally, with varying degrees of sophistication and
diplomatic effect. Table~\ref{tab:systems} provides a comparative
overview of the most relevant systems.

\begin{table}[ht]
\caption{Comparative Digital Travel Declaration Systems}
\label{tab:systems}
\small
\begin{tabularx}{\columnwidth}{@{}lXll@{}}
\toprule
\textbf{System} & \textbf{Function} & \textbf{Year} & \textbf{Outcome} \\
\midrule
US ESTA         & Pre-travel authorisation for visa-waiver travellers & 2009 & Risk stratification \\
Canada eTA      & Electronic travel authorisation, linked to biometrics & 2015 & Reduced fraud \\
Australia ETA   & Electronic travel authority, real-time processing & 1996 & Model for APAC \\
Philippines eTravel & Pre-departure/arrival declaration for all travellers & 2022 & Canada visa-free 2024 \\
EU EES          & Entry/exit system tracking Schengen overstays & 2024 & Overstay monitoring \\
EU ETIAS        & Pre-travel screening for visa-exempt travellers & 2025 & Risk assessment \\
\bottomrule
\end{tabularx}
\end{table}

\subsection{The Philippines eTravel System: Primary Reference}

The Philippines' eTravel system, launched in 2022 under the Bureau of
Immigration and the Department of Information and Communications
Technology, is the most directly analogous reference system for the
proposed NDTR. It is mandatory for all travellers entering or departing
the Philippines---including foreign nationals---and requires
pre-registration of trip details, accommodation, and health declarations
through a mobile application or web portal~\cite{philippinesetravel2022}.

The system was personally used by the author during a visit to the Philippines in
January 2026. The eTravel platform---accessed via the \textit{eGovPH}
super-application---generated unique QR code-based travel declarations
for both arrival (Ireland to Philippines, 14 January 2026, cleared at
Ninoy Aquino International Airport Terminal 3) and departure
(Philippines to Ireland, 20 January 2026). The system maintains a
timestamped, verifiable travel history accessible to the traveller and
to immigration officers.

The key governance value of the Philippines eTravel system lies not in
its UX but in its data architecture: every trip generates a unique
reference number linked to passport details, declared purpose,
accommodation address, and flight routing. This creates a nationally
aggregated compliance database. By 2024, the Philippines leveraged its
broader digital governance improvements---of which eTravel is a
component---in bilateral negotiations that resulted in Canada granting
visa-free access to Philippine passport holders~\cite{canada2024ph}.
This is the diplomatic feedback loop that the NDTR seeks to replicate
for Nigeria.

Significantly, the Philippines itself restricts Nigerian passport entry,
requiring a visa. The fact that a Nigerian traveller is required to
use the Philippines' digital travel declaration system to enter that
country---while Nigeria operates no equivalent system for its own
outbound citizens---illustrates precisely the asymmetry the NDTR is
designed to address.

\subsection{Limitations of Existing Systems for Nigeria}

Existing systems in the Nigerian context include the Nigeria Immigration
Service's e-passport and the ECOWAS travel certificate. However, neither
constitutes a \textit{behavioural data registry}: they are identity
documents rather than trip-level declaration systems. The critical
distinction is that the NDTR is not about travel authorisation
(which destination countries control) but about \textit{travel declaration}
by Nigerian citizens to their own government---creating a national data
asset from outbound mobility patterns.

%% ----------------------------------------------------------
\section{The NDTR Framework}
\label{sec:framework}
%% ----------------------------------------------------------

\subsection{Design Principles}

The NDTR is designed around five core principles:

\begin{enumerate}[label=\textbf{P\arabic*.}, leftmargin=*]
  \item \textbf{Mandatory universality.} All Nigerian citizens travelling
  internationally must register each trip, both outbound and inbound.
  No exceptions below diplomatic exemption level.

  \item \textbf{NIN anchoring.} Each travel declaration is linked to the
  traveller's National Identity Number, creating a verified, persistent
  identity anchor across all declarations.

  \item \textbf{Minimal friction design.} The registration interface must
  be completable in under 5 minutes via mobile application, USSD, or
  web portal, to ensure accessibility across the Nigerian digital divide.

  \item \textbf{Data sovereignty.} All raw declaration data resides within
  Nigerian jurisdiction, under the governance of the National Data
  Protection Commission (NDPC), with strict controls on external access.

  \item \textbf{Structured openness.} Aggregated, anonymised compliance
  statistics are made available to bilateral treaty partners and
  international organisations under formal data-sharing agreements.
\end{enumerate}

\subsection{Core Data Fields}

Each NDTR declaration captures the following minimum dataset:

\begin{itemize}
  \item Traveller NIN and passport number
  \item Origin and destination country
  \item Declared purpose of travel (education, business, tourism, medical,
  family, other)
  \item Planned duration of stay
  \item Accommodation address (destination)
  \item Carrier and flight/vessel identifier
  \item Emergency contact in Nigeria
  \item Return declaration (confirmed upon re-entry)
\end{itemize}

The return declaration is critical: it closes the loop on each trip,
generating a \textit{compliance event} when a traveller returns within
their declared duration. Trips without a return declaration within a
configurable threshold period generate an alert for follow-up by NIS.

\subsection{Technical Architecture}

The NDTR operates as a central registry managed by the Nigeria
Immigration Service, with the following technical layers:

\begin{description}
  \item[Identity Layer] NIN validation via the National Identity
  Management Commission (NIMC) API, ensuring that each declarant is a
  verified Nigerian citizen.

  \item[Declaration Layer] Mobile/web/USSD interface for trip
  declaration, generating a unique QR reference code per trip (similar
  to the Philippines eTravel QR code), required for check-in at
  Nigerian international airports.

  \item[Analytics Layer] Aggregate compliance dashboards for NIS
  policymakers, producing national-level statistics on travel purpose,
  destination, duration, and return compliance rates.

  \item[Security Layer] Anomaly detection module (detailed in
  Section~\ref{sec:security}) that flags atypical patterns for
  targeted review.

  \item[Governance Layer] Data retention policies, audit trails, and
  NDPC-compliant access control, ensuring that individual data is
  not accessible without lawful basis.
\end{description}

\subsection{Privacy and Data Governance}

Privacy by design is non-negotiable for NDTR legitimacy and public
trust. The framework proposes the following safeguards:

\begin{itemize}
  \item Traveller-facing data transparency: citizens can view, contest,
  and request deletion of their own records via a self-service portal.

  \item No commercial data use: NDTR data may only be used for
  immigration governance, national security, and bilateral treaty purposes.

  \item Tiered access: raw individual records accessible only to NIS
  with judicial oversight; aggregated statistics accessible more broadly.

  \item Bilateral data-sharing only under treaty: no unilateral data
  transfer to foreign governments without explicit bilateral agreement
  and NDPC approval.

  \item Sunset provisions: individual trip records deleted after
  10 years unless flagged by an active security review.
\end{itemize}

%% ----------------------------------------------------------
\section{Proactive Security Monitoring and Anomaly Detection}
\label{sec:security}
%% ----------------------------------------------------------

A critical and underexplored dimension of digital travel registries
is their capacity for \textit{proactive} rather than reactive security
monitoring. Existing Nigerian border security practices are largely
reactive: action is taken after an incident is reported by a foreign
authority. The NDTR enables a fundamentally different operational mode.

\subsection{Behavioural Baseline and Anomaly Detection}

As the NDTR accumulates trip-level data, it builds a national
behavioural baseline: the typical frequency, destination, duration,
and purpose distribution for Nigerian international travel. Anomalies
relative to this baseline---at either the individual or pattern
level---become detectable without targeting any specific individual
a priori.

Examples of detectable anomaly classes include:

\begin{description}
  \item[Destination clustering] An unusual concentration of travellers
  declaring the same destination city within a short window, with
  no corresponding event (conference, festival, sporting event), may
  indicate organised irregular migration activity.

  \item[Purpose-duration mismatch] A traveller who declares a
  2-day business trip to a country with no significant Nigerian
  commercial presence presents a pattern worthy of secondary review.

  \item[Rapid repeat travel] Multiple short-duration trips to the
  same destination within a 30-day window, particularly to known
  transit hubs, may indicate courier activity.

  \item[Non-return alerts] Travellers who do not submit a return
  declaration within a threshold period beyond their declared
  stay trigger an automated alert, allowing NIS to initiate
  welfare or enforcement inquiries.

  \item[High-risk corridor flagging] Declared travel to known
  drug-transit or conflict-zone corridors can trigger additional
  pre-departure checks in coordination with the National Drug Law
  Enforcement Agency (NDLEA) and the Department of State Services
  (DSS).
\end{description}

\subsection{Incident Prevention vs.\ Incident Response}

The distinction between prevention and response is operationally
significant.
Under current arrangements, Nigeria learns of an incident involving
one of its citizens abroad---a drug arrest, an immigration
violation, a security incident---through a foreign government
notification, typically after the fact. The NDTR reframes this
operationally: NIS has visibility of anomalous outbound travel
patterns in real time, and can share relevant intelligence with
partner agencies or destination-country counterparts before
departure or upon arrival.

This transforms the NDTR from a passive registry into an active
component of Nigeria's national security infrastructure---an
argument that is likely to be decisive in securing interagency
buy-in from NIS, NDLEA, the DSS, and the Office of the
National Security Adviser.

\subsection{International Security Cooperation}

The NDTR's security monitoring layer is designed to be interoperable
with INTERPOL's I-24/7 communications network and the ICAO
traveller identification programme, enabling Nigeria to share
flagged travel intelligence with foreign partners on a
reciprocal basis---a capability that currently does not exist
at any meaningful scale~\cite{icao2023}.

%% ----------------------------------------------------------
\section{Diplomatic Applications of NDTR Data}
\label{sec:diplomatic}
%% ----------------------------------------------------------

\subsection{Countering Visa Restrictions with Compliance Evidence}

The primary diplomatic application of the NDTR is the generation
of \textit{aggregate compliance evidence}---a national-level
dataset demonstrating that the overwhelming majority of Nigerian
travellers honour their declared purpose and return within the
stated duration.

With three to five years of NDTR data, the Ministry of Foreign Affairs
gains a qualitatively new instrument for bilateral visa negotiations.
Instead of asserting Nigerian traveller compliance in the absence of
evidence, it can table a dataset: in a given year, X million Nigerians
travelled to Y destination; Z\% declared a purpose of business or
education; W\% returned within their declared duration;
anomaly-flagged cases represented P\% of total travel.
This is the same type of argument that gave the Philippines leverage
in its negotiations with Canada.

For ECOWAS and African Union free movement negotiations, aggregate
compliance data is a direct strengthening factor.
For Schengen negotiations, the European Union has published specific
criteria for bilateral visa liberalisation dialogues that include
data on irregular migration rates and document security---criteria
for which NDTR data is directly relevant~\cite{euvisadialogue2023}.

\subsection{Asylum Claim Integrity}

As discussed in Section~\ref{sec:background}, the NDTR provides
a mechanism for asylum claim integrity verification.
Under a bilateral data-sharing arrangement with a destination country's
immigration authority, that authority can---with the claimant's consent,
or under a mutual legal assistance treaty---request an individual's NDTR
record. A claimant's history of regular, undisturbed international
travel, particularly to countries with no relevant asylum context, is
legally and logically relevant to the credibility of a persecution claim.

The NDTR does not determine the outcome of an asylum claim.
It provides verifiable, official Nigerian government data to
an adjudication process that currently operates with incomplete
information. This benefits genuine claimants and deters fraudulent ones.

\subsection{Credit and Financial Obligation Compliance}

A secondary application is the potential integration of NDTR data
with financial compliance monitoring. Nigerian citizens who travel
internationally while leaving unresolved financial obligations---unpaid
loans, tax arrears, outstanding court orders---represent a risk category
for which cross-border visibility is currently near-zero.
With NIN integration, the NDTR provides a technical basis (under
appropriate legal framework) for coordinated notification between
NIS, the Central Bank of Nigeria (CBN), and relevant financial
institutions---analogous to mechanisms that operate in several
EU member states~\cite{cbn2023}. This application requires
separate legislative authority and is flagged as a second-phase
consideration.

%% ----------------------------------------------------------
\section{Border Management Modernisation}
\label{sec:border}
%% ----------------------------------------------------------

\subsection{Alignment with ICAO Standards}

The International Civil Aviation Organization's Traveller Identification
Programme (TRIP) provides the international standards framework within
which the NDTR must operate. ICAO Doc~9303 governs machine-readable
travel documents, and ICAO's TRIP strategy promotes the linkage of
identity management, pre-departure data sharing, and border control
into a coherent national system~\cite{icao2023}. The NDTR is designed
to be ICAO-compliant at the data layer, ensuring that Nigerian-generated
travel data can be exchanged in ICAO-standard formats with partner
state border agencies.

\subsection{ECOWAS Free Movement and the African Union Protocol}

The Economic Community of West African States (ECOWAS) Free Movement
Protocol theoretically permits citizens of member states to travel
freely within the region without a visa for up to 90 days.
In practice, Nigerian travellers continue to face administrative
barriers at land borders across the ECOWAS zone, and the protocol
lacks a digital implementation layer.
The NDTR provides a mechanism for Nigeria to lead the
operationalisation of ECOWAS free movement through digital means:
by demonstrating national-level compliance infrastructure,
Nigeria can advocate for a common ECOWAS digital travel declaration
standard modelled on the NDTR~\cite{ecowas2022}.

At the continental level, the African Union's Free Movement
of Persons Protocol envisages a continental digital passport and
identity infrastructure. The NDTR, built on NIN, positions Nigeria's
migration data infrastructure for interoperability with the AU's
eventual continental digital governance architecture.

\subsection{Modernising Nigerian Airport Operations}

The operational integration of NDTR QR codes into the check-in
and immigration clearance workflow at Nigerian international
airports---beginning with Murtala Muhammed International Airport
(Lagos), Nnamdi Azikiwe International Airport (Abuja), and
Mallam Aminu Kano International Airport (Kano)---would provide
the physical point-of-use infrastructure for the system.
The QR-based check-in model is directly analogous to the
Philippines eTravel QR code, which is presented at the immigration
counter rather than requiring officer data entry, reducing
queuing times and eliminating manual transcription errors.

%% ----------------------------------------------------------
\section{Implementation Roadmap}
\label{sec:roadmap}
%% ----------------------------------------------------------

The NDTR is proposed as a phased implementation over 36 months,
recognising the institutional complexity of interagency coordination
in the Nigerian public sector.

\subsection{Phase 1: Foundation (Months 1--12)}

\begin{itemize}
  \item Legislative basis: amendment to the Nigeria Immigration Act to
  mandate digital travel declarations for outbound citizens.
  \item Interagency governance structure: establish a cross-ministry
  NDTR Steering Committee chaired by NITDA, with representation from
  NIS, NIMC, NDPC, the Ministry of Foreign Affairs, and the NSA office.
  \item Core NIN-NDTR integration: technical design and procurement for
  the declaration platform, built on the existing NIMC NIN API.
  \item Pilot deployment: Lagos (MMIA) outbound terminal for a
  6-month voluntary registration pilot, generating baseline data
  and UX feedback.
\end{itemize}

\subsection{Phase 2: Rollout (Months 13--24)}

\begin{itemize}
  \item Mandatory declaration enforcement at all five international
  airports (Lagos, Abuja, Kano, Port Harcourt, Enugu).
  \item Anomaly detection module deployment on the analytics layer.
  \item First annual compliance report: NDTR publishes aggregated
  national travel statistics for the preceding 12-month period.
  \item Initial bilateral engagement: Nigeria tables NDTR data in
  visa dialogue with at least two partner states (proposed: UK and
  Netherlands as Schengen representative).
\end{itemize}

\subsection{Phase 3: Diplomatic Operationalisation (Months 25--36)}

\begin{itemize}
  \item Land border integration: extend NDTR declaration requirement
  to major land border crossings (beginning with the Seme and
  Idiroko crossings on the Lagos axis).
  \item ECOWAS pilot: propose a common ECOWAS digital travel
  declaration standard based on NDTR architecture.
  \item Bilateral data-sharing agreements: conclude formal MoUs
  with at least three destination countries for reciprocal
  compliance data exchange.
  \item Asylum integrity protocol: activate the bilateral asylum
  claim verification mechanism with Ireland and the United Kingdom
  under the mutual legal assistance framework.
\end{itemize}

\subsection{Key Stakeholders and Roles}

\begin{table}[ht]
\caption{NDTR Stakeholder Roles}
\label{tab:stakeholders}
\small
\begin{tabularx}{\columnwidth}{@{}lX@{}}
\toprule
\textbf{Agency} & \textbf{Role} \\
\midrule
NITDA                & Technical lead; platform development and programme oversight \\
NIS                  & Operational lead; border enforcement integration \\
NIMC                 & NIN API and identity verification \\
NDPC                 & Data governance, privacy compliance, audit \\
Min.\ Foreign Affairs & Bilateral diplomatic deployment of NDTR compliance data \\
NDLEA / DSS          & Security anomaly review and inter-agency alert handling \\
CBN                  & Phase 2: financial obligation integration (requires separate legislation) \\
\bottomrule
\end{tabularx}
\end{table}

%% ----------------------------------------------------------
\section{Discussion}
\label{sec:discussion}
%% ----------------------------------------------------------

The NDTR framework addresses a governance gap that is simultaneously
technical, diplomatic, and philosophical. Technically, it fills the
absence of a national behavioural travel database. Diplomatically,
it converts that database into bilateral leverage. Philosophically,
it reframes the terms of Nigeria's engagement with the global migration
system---from a reactive posture of defending Nigerian citizens against
collective stigma, to a proactive one of demonstrating, through evidence,
why that stigma is statistically unjustified.

There are legitimate concerns to address. First, the risk of government
surveillance overreach: a mandatory travel registry, if misused, could
become an instrument of political control over citizens' freedom of
movement. The data governance framework proposed in
Section~\ref{sec:framework} is designed precisely to prevent this, and
the NDPC's independence is a critical institutional safeguard.

Second, the risk of exclusion: Nigerians without smartphones or reliable
internet access---a non-trivial proportion of the travel-eligible
population, particularly for land border crossings---must not be
systemically excluded. The USSD channel and airport kiosk options are
essential, not optional.

Third, the risk of data breach: a centralised national travel registry
is a high-value target. The security architecture must treat breach
prevention as a first-order design requirement, not an afterthought.

These concerns are manageable and do not negate the fundamental case
for the NDTR. They do require that the implementation be designed with
civil society oversight, independent audit, and a clear legislative
mandate rather than administrative decree.

%% ----------------------------------------------------------
\section{Policy Recommendations}
\label{sec:recommendations}
%% ----------------------------------------------------------

Based on the framework developed in this document, the following
policy recommendations are directed to NITDA and partner agencies
for consideration:

\begin{enumerate}[label=\textbf{R\arabic*.}, leftmargin=*, itemsep=6pt]

  \item \textbf{Establish an NDTR Interagency Steering Committee.}
  NITDA should convene and chair a formal interagency body comprising
  NIS, NIMC, NDPC, the Ministry of Foreign Affairs, and the Office of
  the National Security Adviser. This body should be mandated with
  producing a feasibility report within 6 months of formation.

  \item \textbf{Commission a legal gap analysis.}
  Identify the specific legislative amendments required to the Nigeria
  Immigration Act and the NIN Administration Act to provide a mandatory
  digital declaration requirement with appropriate enforcement provisions
  and data protection safeguards.

  \item \textbf{Procure a voluntary pilot at MMIA Lagos.}
  Deploy a 6-month voluntary NDTR registration pilot at the Murtala
  Muhammed International Airport outbound terminal, integrated with the
  NIMC NIN API. The pilot should target at minimum 100,000 declarations,
  generating the dataset needed to validate the architecture and
  demonstrate viability.

  \item \textbf{Publish an annual NDTR Compliance Report.}
  Within 24 months of full launch, NITDA and NIS should co-publish an
  annual national travel compliance report, providing aggregate statistics
  on travel purpose, return rates, destination distribution, and
  anomaly-flagged cases. This report is the primary instrument for
  diplomatic use.

  \item \textbf{Initiate bilateral data diplomacy.}
  The Ministry of Foreign Affairs should be briefed on the NDTR data
  framework and tasked with initiating formal data-sharing dialogue with
  the United Kingdom Home Office and the Dutch Immigration and
  Naturalisation Service (IND) within the first 36 months of NDTR
  operation, using the first annual compliance report as the opening
  instrument.

  \item \textbf{Prioritise USSD and kiosk access channels.}
  The procurement specification for the NDTR platform must include USSD
  and airport self-service kiosk channels as mandatory requirements,
  not optional enhancements, to prevent a digital-divide exclusion
  problem that would undermine the system's mandatory universality
  principle.

  \item \textbf{Engage ECOWAS on a regional digital declaration standard.}
  Nigeria should use the NDTR architecture as the basis for a regional
  proposal to the ECOWAS Commission for a common digital travel
  declaration standard across ECOWAS member states, positioning Nigeria
  as the technical lead and accelerating the operationalisation of the
  ECOWAS Free Movement Protocol.

\end{enumerate}

%% ----------------------------------------------------------
\section{Conclusion}
\label{sec:conclusion}
%% ----------------------------------------------------------

This document has presented the Nigerian Digital Travel Registry (NDTR)
as a policy instrument that addresses four interconnected governance
deficits: the inability to demonstrate traveller compliance to foreign
governments; the absence of proactive security monitoring for atypical
outbound travel patterns; the lack of verifiable data for Nigeria's
bilateral visa negotiations; and the gap between Nigerian border
management practices and international digital governance standards.

The Philippines eTravel system demonstrates that a digital travel
declaration framework, implemented at scale, can genuinely shift
a country's position in the global migration system.
Nigeria has the demographic scale, the NIN identity infrastructure,
the technical talent embedded in NITDA, and the political impetus---
through the Ministry of Communications and Digital Economy's
broader digital transformation agenda---to replicate and exceed
that model.

The NDTR is not a solution to all of Nigeria's passport restrictions.
Geopolitical factors, trade relationships, and historical migration
patterns all play a role in bilateral visa policy.
But it addresses the one factor that Nigerian governance can
directly control: the availability of credible, government-generated
data about the travel behaviour of its own citizens.
In the absence of such data, Nigeria will continue to negotiate
with one hand tied behind its back.
The NDTR is the instrument that removes that constraint.

This document is presented as a working framework intended to invite
engagement, critique, and collaborative development with NITDA and
relevant agencies. The author welcomes further dialogue and is
available to elaborate on any aspect of the proposed framework.

\vspace{8pt}
\paragraph{Contact.}
Olajide Adewale Olaniyan \quad\textbar\quad
\href{mailto:jidemobell@gmail.com}{jidemobell@gmail.com} \quad\textbar\quad
IBM Ireland, Dublin, Ireland

\paragraph{Acknowledgements.}
The author acknowledges the initial reception of this proposal
by the National Information Technology Development Agency (NITDA),
whose engagement in January 2026 encouraged the development of this
framework document.

%% ----------------------------------------------------------
%%  REFERENCES
%% ----------------------------------------------------------

\bibliographystyle{unsrt}

\begin{thebibliography}{99}

\bibitem{henley2025}
Henley \& Partners.
\newblock \textit{Henley Passport Index 2025}.
\newblock Henley \& Partners, 2025.
\newblock \url{https://www.henleyglobal.com/passport-index}

\bibitem{iom2023}
International Organization for Migration (IOM).
\newblock \textit{World Migration Report 2024}.
\newblock IOM, Geneva, 2023.

\bibitem{unodc2022}
United Nations Office on Drugs and Crime (UNODC).
\newblock \textit{Global Report on Trafficking in Persons 2022}.
\newblock UNODC, Vienna, 2022.

\bibitem{castles2013}
Stephen Castles, Hein De Haas, and Mark J. Miller.
\newblock \textit{The Age of Migration: International Population Movements
in the Modern World}, 5th ed.
\newblock Palgrave Macmillan, 2013.

\bibitem{kahneman2011}
Daniel Kahneman.
\newblock \textit{Thinking, Fast and Slow}.
\newblock Farrar, Straus and Giroux, 2011.

\bibitem{eurostat2024}
Eurostat.
\newblock \textit{Asylum and Managed Migration Statistics 2023}.
\newblock European Commission, 2024.
\newblock \url{https://ec.europa.eu/eurostat/statistics-explained/index.php/Asylum_statistics}

\bibitem{irishnat2023}
International Protection Appeals Tribunal (Ireland).
\newblock \textit{Annual Report 2023}.
\newblock Department of Justice, Ireland, 2023.

\bibitem{philippinesetravel2022}
Department of Information and Communications Technology (DICT), Philippines.
\newblock \textit{eTravel System: Technical and Operational Overview}.
\newblock DICT, Quezon City, 2022.
\newblock \url{https://etravel.gov.ph}

\bibitem{canada2024ph}
Immigration, Refugees and Citizenship Canada (IRCC).
\newblock \textit{Visa-Free Travel for Philippine Passport Holders}, 2024.
\newblock \url{https://www.canada.ca}

\bibitem{euvisadialogue2023}
European Commission.
\newblock \textit{Visa Liberalisation Action Plans: Monitoring and Evaluation
Framework}.
\newblock Directorate-General for Migration and Home Affairs, Brussels, 2023.

\bibitem{cbn2023}
Central Bank of Nigeria (CBN).
\newblock \textit{Framework for Cross-Border Financial Intelligence Sharing}.
\newblock CBN, Abuja, 2023.

\bibitem{icao2023}
International Civil Aviation Organization (ICAO).
\newblock \textit{ICAO Traveller Identification Programme (TRIP) Strategy:
Third Edition}.
\newblock ICAO, Montreal, 2023.
\newblock Doc 9303.

\bibitem{ecowas2022}
Economic Community of West African States (ECOWAS).
\newblock \textit{Protocol A/P.1/5/79 on Free Movement of Persons:
2022 Progress Report}.
\newblock ECOWAS Commission, Abuja, 2022.

\end{thebibliography}

\end{document}
